% ============================================================
% Petites macros TikZ pour dessiner les diagrammes UML du rapport
% (acteurs, classes, séquence) sans aucun outil externe.
% ============================================================

\definecolor{umlgrey}{gray}{0.55}
\definecolor{umlfill}{gray}{0.93}

\tikzset{
  >=Stealth,
  uml/.style={font=\footnotesize\sffamily},
  % ---------- cas d'utilisation ----------
  uc/.style={ellipse, draw=black, fill=white, align=center, font=\scriptsize\sffamily,
             minimum height=0.95cm, text width=2.9cm, inner sep=1pt},
  ucint/.style={uc, fill=umlfill},
  assoc/.style={draw=black},
  incl/.style={draw=black, dashed, -{Stealth[length=2mm]}},
  gener/.style={draw=black, -{Triangle[open, length=3mm, width=3mm]}},
  rellabel/.style={font=\fontsize{6}{7}\selectfont\sffamily, fill=white, inner sep=1pt, sloped},
  % ---------- classes ----------
  cls/.style={rectangle split, rectangle split parts=3, draw=black, fill=white,
              rectangle split part align={center,left,left}, font=\scriptsize\sffamily,
              inner xsep=5pt, inner ysep=3pt, text width=3.7cm},
  clsenum/.style={cls, rectangle split parts=2, rectangle split part align={center,left}, fill=umlfill},
  mult/.style={font=\fontsize{6}{7}\selectfont\sffamily, inner sep=1.5pt},
  % ---------- séquences ----------
  lifebox/.style={draw=black, fill=white, align=center,
                  font=\scriptsize\sffamily, inner sep=3pt, minimum height=0.7cm, text width=2.5cm},
  lifeline/.style={draw=black!60, dashed, thin},
  msg/.style={-{Stealth[length=2mm]}, draw=black},
  retmsg/.style={-{Stealth[length=2mm, open]}, dashed, draw=black},
  msgl/.style={above, font=\scriptsize\sffamily, align=center, inner sep=1pt, fill=white, fill opacity=0.85, text opacity=1},
  fragbox/.style={draw=black!60},
  fragtab/.style={draw=black!60, fill=white, font=\scriptsize\sffamily, inner sep=2pt, anchor=north west},
  fragcond/.style={font=\scriptsize\sffamily, anchor=west},
}

% ---------- acteur (bonhomme) : \actor{nom}{x,y}{libellé} ----------
\newcommand{\actor}[3]{%
  \begin{scope}[shift={(#2)}]
    \draw (0,0.62) circle (0.17);
    \draw (0,0.45)--(0,-0.05);
    \draw (-0.3,0.3)--(0.3,0.3);
    \draw (0,-0.05)--(-0.22,-0.5);
    \draw (0,-0.05)--(0.22,-0.5);
  \end{scope}
  \node[inner sep=0, minimum width=0.7cm, minimum height=1.1cm] (#1) at ($(#2)+(0,0.05)$) {};
  \node[font=\scriptsize\sffamily\bfseries, align=center, below=0.5cm of #1, text width=2.4cm] {#3};
}

% ---------- séquences ----------
% Les lignes de vie sont désignées par un identifiant ; \lx{id} donne leur abscisse (en cm).
\newcommand{\lx}[1]{\csname lx@#1\endcsname}
\newcommand{\seqbegin}{\def\cury{-0.35}\gdef\lifelist{}\def\seqk{0.85}}
\newcommand{\life}[3]{%
  \expandafter\gdef\csname lx@#1\endcsname{#2}%
  \node[lifebox] at (#2,0) {#3};
  \ifx\lifelist\empty\xdef\lifelist{#2}\else\xdef\lifelist{\lifelist,#2}\fi
}
\newcommand{\seqend}{%
  \pgfmathsetmacro{\yend}{\cury-0.5}%
  \begin{scope}[on background layer]
    \foreach \lfx in \lifelist {\draw[lifeline] (\lfx,-0.35)--(\lfx,\yend);}
  \end{scope}
}
% message : \msg[pas]{de}{vers}{libellé}   (\ret : retour en pointillé)
\newcommand{\msg}[4][0.8]{%
  \pgfmathsetmacro{\cury}{\cury-#1*\seqk}%
  \pgfmathsetmacro{\mxa}{\lx{#2}}\pgfmathsetmacro{\mxb}{\lx{#3}}%
  \draw[msg] (\mxa,\cury)--(\mxb,\cury) node[msgl, pos=0.5] {#4};%
}
\newcommand{\ret}[4][0.8]{%
  \pgfmathsetmacro{\cury}{\cury-#1*\seqk}%
  \pgfmathsetmacro{\mxa}{\lx{#2}}\pgfmathsetmacro{\mxb}{\lx{#3}}%
  \draw[retmsg] (\mxa,\cury)--(\mxb,\cury) node[msgl, pos=0.5] {#4};%
}
% message sur soi-même : \selfmsg[pas]{ligne}{libellé}
\newcommand{\selfmsg}[3][0.8]{%
  \pgfmathsetmacro{\cury}{\cury-#1*\seqk}%
  \pgfmathsetmacro{\sx}{\lx{#2}}%
  \pgfmathsetmacro{\ya}{\cury-0.38}%
  \pgfmathsetmacro{\xb}{\sx+0.55}%
  \pgfmathsetmacro{\xc}{\sx+0.6}%
  \pgfmathsetmacro{\ym}{\cury-0.19}%
  \draw[msg] (\sx,\cury)--(\xb,\cury)--(\xb,\ya)--(\sx,\ya);%
  \node[font=\scriptsize\sffamily, anchor=west, align=left, inner sep=1.5pt, fill=white, fill opacity=0.85, text opacity=1] at (\xc,\ym) {#3};%
  \pgfmathsetmacro{\cury}{\cury-0.38}%
}
% fragment combiné : \fragbegin[marge]{id}{type}{ligne gauche}{ligne droite}{condition}
%                    \fragelse{id}{condition}   \fragend{id}
\newcommand{\fragbegin}[6][0.7]{%
  \pgfmathsetmacro{\cury}{\cury-0.3}%
  \pgfmathsetmacro{\fxl}{\lx{#4}-#1}\pgfmathsetmacro{\fxr}{\lx{#5}+#1}%
  \expandafter\xdef\csname fy@#2\endcsname{\cury}%
  \expandafter\xdef\csname fx1@#2\endcsname{\fxl}%
  \expandafter\xdef\csname fx2@#2\endcsname{\fxr}%
  \node[fragtab] (tab#2) at (\fxl,\cury) {#3};%
  \node[fragcond] at ($(tab#2.east)+(0.1,0)$) {#6};%
  \pgfmathsetmacro{\cury}{\cury-0.3}%
}
\newcommand{\fragelse}[2]{%
  \pgfmathsetmacro{\cury}{\cury-0.4}%
  \expandafter\let\expandafter\fxa\csname fx1@#1\endcsname
  \expandafter\let\expandafter\fxb\csname fx2@#1\endcsname
  \draw[fragbox, dashed] (\fxa,\cury)--(\fxb,\cury);%
  \pgfmathsetmacro{\xe}{\fxa+0.1}%
  \pgfmathsetmacro{\ye}{\cury-0.2}%
  \node[fragcond] at (\xe,\ye) {#2};%
  \pgfmathsetmacro{\cury}{\cury-0.25}%
}
\newcommand{\fragend}[1]{%
  \pgfmathsetmacro{\yfe}{\cury-0.3}%
  \expandafter\let\expandafter\fxa\csname fx1@#1\endcsname
  \expandafter\let\expandafter\fxb\csname fx2@#1\endcsname
  \expandafter\let\expandafter\fya\csname fy@#1\endcsname
  \draw[fragbox] (\fxa,\fya) rectangle (\fxb,\yfe);%
  \pgfmathsetmacro{\cury}{\cury-0.15}%
}

% ---------- classes : \umlclass[options]{nom}{style}{titre}{attributs}{méthodes} ----------
\newcommand{\umlclass}[6][]{%
  \node[#3, #1] (#2) {#4\nodepart{second}#5\nodepart{third}#6};%
}
\newcommand{\umlenum}[4][]{%
  \node[clsenum, #1] (#2) {#3\nodepart{second}#4};%
}

% Figure UML mise à l'échelle de la largeur de ligne
\newcommand{\umlfig}[2][0.62]{%
  \begingroup\hyphenpenalty=10000 \exhyphenpenalty=10000
  \noindent\adjustbox{max width=\linewidth, max totalheight=#1\textheight, center}{\input{#2}}%
  \endgroup}
